% 03 Method
\section{方法}
\label{sec:method}

本节详细描述渐进式频率条件 FiLM 框架的整体架构与关键模块。为便于表述，设输入序列为 \(X \in \mathbb{R}^{B \times L \times C}\)，其中 \(B\) 为 batch 大小，\(L\) 为历史序列长度（回望窗口），\(C\) 为变量维度（通道数）。模型整体遵循 iTransformer~\cite{itransformer} 的变量维 token 化范式，将每个变量的长度为 \(L\) 的历史片段映射为一个 \(d_{\text{model}}\) 维 token，并通过 Transformer Encoder 建模跨变量交互。本文的核心贡献在于：(1) 在变量 token 流入 Encoder 之前及每一层之前，注入由 FFT 多分辨率频谱生成的条件调制信号；(2) 调制采用逐层独立生成、强度随深度指数衰减的渐进式 FiLM 机制；(3) 全局调制基强度为可学习参数，支持数据集自适应与可退化训练。

\subsection{总体架构}
\label{subsec:overview}

渐进式频率条件 FiLM 的前向计算流程如下（见图\ref{fig:architecture}，待补充）：

\begin{enumerate}[leftmargin=*,label=\textbf{Step \arabic*:},itemsep=4pt]
  \item \textbf{RevIN 归一化}：对输入 \(X\) 沿时间维做可逆实例归一化，得到 \(\tilde{X} \in \mathbb{R}^{B \times L \times C}\)，缓解非平稳分布漂移；
  \item \textbf{变量 Token 嵌入}：将每个变量的 \(L\) 步序列线性映射为 \(d_{\text{model}}\) 维 token，并通过 LayerNorm，得到 \(T^{(0)} \in \mathbb{R}^{B \times C \times d_{\text{model}}}\)；
  \item \textbf{多分辨率频谱提取}：在三个时间窗口上分别对 \(\tilde{X}\) 做 FFT 并提取频带能量、高低频比与频谱熵，得到多分辨率频谱特征 \(S \in \mathbb{R}^{B \times C \times D_s}\)；
  \item \textbf{频谱压缩}：通过小型 MLP 将高维频谱特征压缩为低维 latent \(Z \in \mathbb{R}^{B \times C \times d_z}\)，防止频域信息过载；
  \item \textbf{逐层渐进式 FiLM 调制 + Transformer 编码}：对 \(\ell = 0, 1, \dots, N-1\) 层，依次执行：
  \begin{enumerate}
    \item 由第 \(\ell\) 层的独立 FiLM Generator 从 \(Z\) 生成调制参数 \((\alpha_\ell, \beta_\ell)\)；
    \item 以强度 \(\gamma_\ell = \gamma_{\text{base}} \cdot \delta^{\ell}\) 对 token 做有界 FiLM 调制；
    \item 经 Pre-LN 与 Transformer Encoder Layer 更新 token 表征，得到 \(T^{(\ell+1)}\)。
  \end{enumerate}
  \item \textbf{预测头与 RevIN 反归一化}：将最终 token 表征经 LayerNorm 与线性映射得到未来 \(H\) 步预测，再通过 RevIN 反归一化恢复到原始尺度。
\end{enumerate}

对于 CEMP 参数回归等 sequence-to-vector 任务，Step 6 中的预测头替换为回归头（token 维池化后接 MLP 映射到参数向量），且根据任务需要可选择是否执行 RevIN 反归一化。

\subsection{RevIN：可逆实例归一化}
\label{subsec:revin}

RevIN~\cite{revin} 对每个样本沿时间维独立计算均值与标准差，将序列归一化到零均值单位方差的分布，并在输出端通过反归一化恢复原始尺度。其形式为：

\begin{align}
  \tilde{X} &= \gamma \odot \frac{X - \mu}{\sigma + \epsilon} + \beta, \quad
  \mu = \frac{1}{L}\sum_{t=1}^{L} X_{:,t,:},\;
  \sigma = \sqrt{\frac{1}{L}\sum_{t=1}^{L} (X_{:,t,:} - \mu)^2 + \epsilon}, \\
  \hat{Y} &= \frac{\hat{Y}_{\text{raw}} - \beta}{\gamma} \odot \sigma + \mu,
\end{align}

其中 \(\gamma, \beta \in \mathbb{R}^{C}\) 为可学习的仿射参数，\(\mu, \sigma\) 在归一化时 detach（不参与梯度回传）。通过 RevIN，模型的骨干部分始终在统计量对齐的分布上学习，有效降低了非平稳性对训练与泛化的干扰。

\subsection{变量 Token 嵌入}
\label{subsec:token_embed}

遵循 iTransformer 的设计，将每个变量的整个历史序列片段作为一个 token：
\begin{equation}
  T^{(0)} = \mathrm{LayerNorm}\big(\phi(\tilde{X}^{\mathsf{T}})\big),
  \quad T^{(0)} \in \mathbb{R}^{B \times C \times d_{\text{model}}},
\end{equation}
其中 \(\phi: \mathbb{R}^{L} \to \mathbb{R}^{d_{\text{model}}}\) 为单层线性映射（无激活函数），\(\tilde{X}^{\mathsf{T}} \in \mathbb{R}^{B \times C \times L}\) 为转置后的归一化序列。该设计将计算复杂度从时间维注意力的 \(O(L^2)\) 降为变量维注意力的 \(O(C^2)\)，尤其适合变量数适中的多变量预测场景。

\subsection{多分辨率频谱特征提取}
\label{subsec:multires_spectrum}

本文提出多分辨率频谱提取器（Multi-Resolution Spectrum Extractor），在三个时间尺度上独立计算 FFT 功率谱并提取结构化的频域特征。该模块的输入为归一化后的序列 \(\tilde{X}^{\mathsf{T}} \in \mathbb{R}^{B \times C \times L}\)，输出为多分辨率频谱特征 \(S \in \mathbb{R}^{B \times C \times D_s}\)。

\medskip\noindent\textbf{多时间窗口.}
在三个递减的时间窗口上分别截取输入序列的末尾部分（最近的信息最重要）：
\begin{equation}
  \mathcal{W} = \{L,\; L/2,\; L/4\},
\end{equation}
对于每个窗口大小 \(w \in \mathcal{W}\)（实际实现中取 \(w \gets \max(8, \lfloor w \rfloor)\)），取 \(X_w = \tilde{X}_{:, :, -w:}^{\mathsf{T}} \in \mathbb{R}^{B \times C \times w}\)。

\medskip\noindent\textbf{功率谱计算.}
对每个窗口的每个变量独立计算实值 FFT（rFFT），并取幅值平方得到功率谱：
\begin{equation}
  \mathcal{F}_w = \mathrm{rFFT}(X_w), \quad
  M_w = |\mathcal{F}_w|^2 \;\in\; \mathbb{R}^{B \times C \times f_w},
\end{equation}
其中频率 bin 数 \(f_w = \lfloor w/2 \rfloor + 1\)。

\medskip\noindent\textbf{每窗口的结构化频域特征.}
对每个窗口的功率谱 \(M_w\)，提取以下三类特征：

\emph{(a) 多频带归一化能量.} 将频率轴均分为 \(K\) 个频带（本文取 \(K = 4\)），计算每频带的平均功率并做 L1 归一化：
\begin{equation}
  E_k^{(w)} = \frac{1}{|\mathcal{B}_k|} \sum_{f \in \mathcal{B}_k} M_w[:,:,f], \quad
  \hat{E}_k^{(w)} = \frac{E_k^{(w)}}{\sum_{j=1}^{K} E_j^{(w)} + \epsilon}, \quad k=1,\dots,K.
\end{equation}
归一化频带能量刻画了信号的功率在不同频率区间的分布形态（例如低频主导 vs 高频主导），且尺度不变。

\emph{(b) 高低频比值.} 取最高频带与最低频带的能量比：
\begin{equation}
  R^{(w)} = \frac{E_K^{(w)}}{E_1^{(w)} + \epsilon}.
\end{equation}
该比值直接反映信号的高频分量相对于低频分量的强弱，是刻画信号"波动性 vs 趋势性"的紧凑归纳偏置。

\emph{(c) 归一化频谱熵.} 将功率谱归一化为概率分布后计算信息熵：
\begin{equation}
  p_f^{(w)} = \frac{M_w[:,:,f]}{\sum_{f'} M_w[:,:,f'] + \epsilon}, \quad
  H^{(w)} = -\frac{\sum_f p_f^{(w)} \log(p_f^{(w)} + \epsilon)}{\log f_w + \epsilon}.
\end{equation}
频谱熵衡量功率在频率轴上的集中程度：低熵对应强周期性/窄带信号（可预测性高），高熵对应白噪声/平坦谱（可预测性低）。除以 \(\log f_w\) 进行归一化使得不同窗口大小的熵值可比。

\medskip\noindent\textbf{特征拼接与数值稳定化.}
将三个窗口的 \(K+2\) 维特征沿最后一维拼接，得到 \(D_s = 3 \times (K+2)\) 维的原始频谱特征：
\begin{equation}
  S_{\text{raw}} = [\hat{E}^{(L)}, R^{(L)}, H^{(L)},\; \hat{E}^{(L/2)}, R^{(L/2)}, H^{(L/2)},\; \hat{E}^{(L/4)}, R^{(L/4)}, H^{(L/4)}].
\end{equation}
随后应用 \(\log(1+x)\) 压缩与截断以保证数值稳定性：
\begin{equation}
  S = \mathrm{clamp}\big(\log(1 + S_{\text{raw}}),\; 0,\; s_{\max}\big), \quad s_{\max}=5.0.
\end{equation}

\subsection{频谱压缩器}
\label{subsec:spectrum_compressor}

为防止高维频谱特征直接注入 Encoder 造成信息过载，我们在频谱特征与 FiLM 生成器之间插入一个频谱压缩器（Spectrum Compressor），将 \(D_s\) 维特征压缩到远小于 \(d_{\text{model}}\) 的低维 latent 空间：
\begin{equation}
  Z = \mathrm{Compressor}(S), \quad
  Z \in \mathbb{R}^{B \times C \times d_z},
\end{equation}
其中 \(\mathrm{Compressor}(\cdot)\) 为一个两层 MLP（带 GELU 激活与 Dropout）：
\begin{equation}
  \mathrm{Compressor}(S) = W_2 \cdot \mathrm{Dropout}\big(\mathrm{GELU}(W_1 \cdot S + b_1)\big) + b_2,
\end{equation}
隐层维度设为 \(\max(64, D_s)\)，latent 维度 \(d_z\) 默认取 32（可配置），显著小于 \(d_{\text{model}}\)（通常为 128--512）。

\subsection{逐层渐进式 FiLM 调制}
\label{subsec:progressive_film}

这是本文最核心的设计——将频域 latent \(Z\) 逐层独立地注入到 Transformer 的每一层之前，并以衰减的强度实现由粗到细的层级化调控。

\medskip\noindent\textbf{可学习调制基强度.}
全局调制强度不由手工设定，而是通过一个可学习的无约束参数 \(\theta_{\gamma}\) 经 Sigmoid 映射得到：
\begin{equation}
  \gamma_{\text{base}} = \sigma(\theta_{\gamma}) = \frac{1}{1 + e^{-\theta_{\gamma}}},
\end{equation}
其中 \(\theta_{\gamma}\) 初始化为 \(\theta_{\gamma}^{(0)} = \ln\frac{0.02}{1-0.02} \approx -3.89\)，使得 \(\gamma_{\text{base}} \approx 0.02\)（保守初始化）。训练过程中，不同数据集自适应学习 \(\theta_{\gamma}\)——若频域信息有效，\(\gamma_{\text{base}}\) 增大；若无效，\(\gamma_{\text{base}}\) 保持较小甚至进一步减小，模型退化为接近纯 iTransformer。

\medskip\noindent\textbf{逐层调制强度衰减.}
对第 \(\ell\) 层（\(\ell = 0, 1, \dots, N-1\)），该层实际使用的调制强度为：
\begin{equation}
  \gamma_{\ell} = \gamma_{\text{base}} \cdot \delta^{\ell},
\end{equation}
其中 \(\delta \in (0, 1)\) 为衰减因子（默认 0.7）。这一设计体现了"由粗到细"的直觉：浅层接受更强的频域引导以建立全局谱结构感知，深层逐渐减弱调制以保护已形成的语义表征不被过度扰动。

\medskip\noindent\textbf{逐层独立 FiLM Generator.}
每一层拥有独立的 FiLM Generator——一个小型 MLP，将共享的频域 latent \(Z\) 映射为该层专属的调制参数：
\begin{equation}
  (\alpha_{\ell}, \beta_{\ell}) = \mathrm{FiLMGen}_{\ell}(Z), \quad
  \alpha_{\ell}, \beta_{\ell} \in \mathbb{R}^{B \times C \times 1}.
\end{equation}

每个 \(\mathrm{FiLMGen}_{\ell}\) 的结构为三层 MLP（含 GELU 激活与 Dropout）：
\begin{equation}
  \mathrm{FiLMGen}(Z) = W_3 \cdot \mathrm{GELU}\big(W_2 \cdot \mathrm{GELU}(W_1 \cdot Z + b_1) + b_2\big) + b_3,
\end{equation}
输出维度为 2（分别对应 \(\alpha\) 和 \(\beta\)），关键设计是最后一层 \(W_3, b_3\) 初始化为零，使得初始状态下 \(\alpha_{\ell} = \beta_{\ell} = 0\)，调制对 token 无影响。

\medskip\noindent\textbf{有界 FiLM 调制.}
在第 \(\ell\) 层 Encoder 之前，对 token 施加经过幅度约束的 FiLM 调制：
\begin{equation}
  T_{\text{mod}}^{(\ell)} = T^{(\ell)} \odot \big(1 + \gamma_{\ell} \cdot \tanh(\alpha_{\ell})\big) + \gamma_{\ell} \cdot \beta_{\ell}.
\end{equation}
其中 \(\tanh(\cdot)\) 将 \(\alpha\) 的缩放效应限制在 \((-1, 1)\) 区间内，乘以小系数 \(\gamma_{\ell}\) 后实际调制幅度在 \((-\gamma_{\ell}, \gamma_{\ell})\) 之间——例如初始状态下约 \((-0.02, 0.02)\)，保证调制仅做"微调"。

\medskip\noindent\textbf{完整逐层计算.}
\begin{equation}
  T^{(\ell+1)} = \mathrm{EncoderLayer}_{\ell}\Big(
    \mathrm{LayerNorm}_{\ell}\big(T_{\text{mod}}^{(\ell)}\big)
  \Big), \quad \ell = 0, 1, \dots, N-1.
\end{equation}
每个 \(\mathrm{EncoderLayer}_{\ell}\) 为标准 Transformer Encoder Layer（Pre-LN 模式，batch\_first=True），包含多头自注意力（变量间交互）与前馈网络。

\subsection{预测头与任务适配}
\label{subsec:head}

经过 \(N\) 层编码后，最终 token 表征经 LayerNorm 与线性映射输出预测：
\begin{equation}
  \hat{Y}_{\text{raw}} = \mathrm{Head}\big(\mathrm{LayerNorm}(T^{(N)})\big),
  \quad \mathrm{Head}: \mathbb{R}^{d_{\text{model}}} \to \mathbb{R}^{H},
\end{equation}
经维度变换后得到 \(\hat{Y}_{\text{raw}} \in \mathbb{R}^{B \times H \times C}\)，再由 RevIN 反归一化恢复到原始尺度。

对于 CEMP 参数回归任务，Head 替换为：Token 维平均池化 \(\to\) 两层 MLP（含 GELU）\(\to\) 参数向量 \(\hat{\bm{p}} \in \mathbb{R}^{P}\)，不执行 RevIN 反归一化。

\subsection{可退化性与稳定化设计总结}
\label{subsec:stability}

渐进式频率条件 FiLM 框架的可退化性由以下三项设计联合保证，三者共同确保模型在初始状态等价于纯 iTransformer，频域调制在训练中"逐步激活"而非"骤然介入"：

\begin{enumerate}[leftmargin=*,itemsep=2pt]
  \item \textbf{FiLM 末层零初始化}：所有 \(\mathrm{FiLMGen}_{\ell}\) 的最后一层 \(W_3, b_3 = 0\)，初始 \(\alpha = \beta = 0\)，调制为恒等操作；
  \item \textbf{保守的 gamma 初始化}：\(\gamma_{\text{base}}\) 初始约 0.02（比常见硬编码值 0.1 更保守），乘以 \(\tanh\) 约束后实际调制幅度极小；
  \item \textbf{数值稳定化管道}：\(\log(1+x)\) 压缩消除长尾分布极端值，clamp 截断防止数值溢出，频谱特征始终在一个可控范围内。
\end{enumerate}

当 \(\gamma_{\text{base}} \to 0\) 时（模型认为频域信息无用），所有层的调制趋近于零，框架严格退化为 iTransformer + RevIN。
